% BHU PYQ -- Manually transcribed & error-corrected
\documentclass[11pt,a4paper]{article}

\usepackage[a4paper,top=2.5cm,bottom=2.5cm,left=2.8cm,right=2.8cm,headheight=14pt]{geometry}
\usepackage{amsmath,amssymb}
\usepackage[T1]{fontenc}
\usepackage[utf8]{inputenc}
\usepackage{lmodern}
\usepackage{microtype}
\usepackage{fancyhdr}
\usepackage{enumitem}
\usepackage{hyperref}
\usepackage{lastpage}
\usepackage{parskip}

\hypersetup{colorlinks=true,urlcolor=black,linkcolor=black,
  pdftitle={BPT-501: Mathematical Physics -- BHU PYQ},pdfauthor={Banaras Hindu University}}

\pagestyle{fancy}\fancyhf{}
\renewcommand{\headrulewidth}{0.4pt}\renewcommand{\footrulewidth}{0.4pt}
\fancyhead[L]{\small   \textit{Banaras Hindu University}}
\fancyhead[R]{\small   \textit{BPT-501: Mathematical Physics}}
\fancyfoot[C]{\small Page  \thepage\ of \pageref{LastPage}}


\newlist{parts}{enumerate}{2}
\setlist[parts,1]{label=(\alph*),leftmargin=*,itemsep=5pt,topsep=3pt}
\setlist[parts,2]{label=(\roman*),leftmargin=*,itemsep=3pt,topsep=2pt}

\newcommand{\pts}[1]{\hfill   \textit{[#1]}}


\begin{document}


\begin{center}
  {\large\bfseries BANARAS HINDU UNIVERSITY}\\[2pt]
  {
\normalsize B.Sc. (Hons.) Semester V Examination, 2022-23}\\[6pt]
  \rule{0.8\linewidth}{0.5pt}\\[6pt]
  {\large\bfseries Paper No. BPT-501: Mathematical Physics}\\[4pt]
  {
\normalsize Time: 3 Hours\hfill Maximum Marks: 70}
\end{center}
\rule{\linewidth}{0.4pt}
\smallskip



\noindent   \textit{Instructions: Answer all questions. Symbols have their usual meanings.}
\bigskip




\noindent   \textbf{Question 1.} Answer any   \textbf{four} of the following: \pts{4 $ \times$ 5 = 20}

\begin{parts}
  \item A conducting wire along the $z$-axis carries a steady current $I$. The resulting magnetic vector potential is given by:
    \[ \vec{A} = \hat{z} \frac{\mu_0 I}{2\pi} \ln \left( \frac{r_0}{r} \right) \]
    Find the magnetic induction vector $\vec{B}$ in coordinates appropriate to the geometry.
  \item The distance square $ds^2$ between two neighboring points in a 3D space is given in curvilinear coordinates by $ds^2 = g_{ij} dx^i dx^j$. Show that the metric coefficients $g_{ij}$ transform as a covariant tensor of rank 2.
  \item If $A$ is a nonsingular symmetric matrix, show that $A^{-1}$ is also symmetric.
  \item From the Euler integral definition of the Gamma function $\Gamma(z) = \int_0^{\infty} e^{-t} t^{z-1} \, dt$, show that:
    
\begin{parts}
      \item $\Gamma(z+1) = z \Gamma(z)$
      \item $\Gamma(z) = 2 \int_0^{\infty} e^{-t^2} t^{2z-1} \, dt$. State the condition imposed on $z$ for the validity of this integral.
    \end{parts}
  \item Find the Fourier series of the function $f(x) = |\sin(x)|$ in the interval $-\pi < x < \pi$. Comment on the nature of convergence of the series.
\end{parts}

\medskip\rule{\linewidth}{0.2pt}




\noindent   \textbf{Question 2.}

\begin{parts}
  \item Normalize the Laguerre polynomial $L_2(x) = 1 - 2x + \frac{x^2}{2}$ in the range $[0, \infty)$ with weight function $e^{-x}$. Construct a second-order, ordinary, linear, homogeneous differential equation (2-ODE) that is satisfied by this polynomial. \pts{6}
  \item Verify the following relations in $\mathbb{R}^3$ using Levi-Civita symbols:
    
\begin{parts}
      \item $\epsilon_{ipq} \epsilon_{ipq} = 6$
      \item $\epsilon_{ijk} \epsilon_{pqk} = \delta_{ip} \delta_{jq} - \delta_{iq} \delta_{jp}$
    \end{parts} \pts{8}
\end{parts}
\hfill   \textbf{OR}

\begin{parts}
  \item Write down the spherical polar unit vectors in terms of Cartesian unit vectors. Using these, calculate all the partial derivatives of $\hat{r}, \hat{ \theta}, \hat{\phi}$ with respect to $r,  \theta, \phi$. \pts{8}
  \item Given a force vector in spherical polar coordinates:
    \[ \vec{F} = \frac{2P \cos \theta}{r^3} \hat{r} + \frac{P \sin \theta}{r^3} \hat{ \theta} \]
    express it in Cartesian coordinates. Does a potential function exist for this force? If so, find it. Provide an example of a physics problem in which the above expression occurs. \pts{6}
\end{parts}

\medskip\rule{\linewidth}{0.2pt}




\noindent   \textbf{Question 3.}

\begin{parts}
  \item Show that the Bessel function $J_{1/2}(x) = \sqrt{\frac{2}{\pi x}} \sin x$ satisfies the Bessel differential equation:
    \[ x^2 y'' + x y' + \left( x^2 - \frac{1}{4} \right) y = 0 \]
    Find a second linearly independent solution of this equation. \pts{7}
  \item Given the generating function $g(x, t) = e^{2xt - t^2} = \sum_{n=0}^{\infty} \frac{H_n(x)}{n!} t^n$ for Hermite polynomials, obtain the Rodrigues representation:
    \[ H_n(x) = (-1)^n e^{x^2} \frac{d^n}{dx^n} \left( e^{-x^2} \right) \] \pts{7}
\end{parts}
\hfill   \textbf{OR}

\begin{parts}
  \item Using the generating function $g(x,t) = (1-2xt+t^2)^{-1/2} = \sum_{n=0}^{\infty} P_n(x)t^n$, $(|t| < 1)$ for Legendre polynomials, show that:
    \[ (n+1)P_{n+1}(x) = (2n+1)x P_n(x) - n P_{n-1}(x) \] \pts{6}
  \item Solve $(1-x^2)y'' - 2xy' + 12y = 0$ by the Frobenius series solution method about $x=0$, and obtain two independent solutions. \pts{8}
\end{parts}

\medskip\rule{\linewidth}{0.2pt}




\noindent   \textbf{Question 4.}

\begin{parts}
  \item By expressing the vector cross products in terms of the Levi-Civita symbols ($\epsilon_{ijk}$), verify the vector identity:
    \[ \vec{A}  \times (\vec{B}  \times \vec{C}) = (\vec{A} \cdot \vec{C})\vec{B} - (\vec{A} \cdot \vec{B})\vec{C} \] \pts{6}
  \item Find the eigenvalues and normalized eigenvectors of the matrix $A$:
    \[ A = 
\begin{pmatrix} 4 & -1 & -1 \\ -1 & 4 & -1 \\ -1 & -1 & 4 \end{pmatrix} \]
    How are these related to the eigenvalues and eigenvectors of $A^2$? \pts{8}
\end{parts}

\medskip\rule{\linewidth}{0.2pt}




\noindent   \textbf{Question 5.}

\begin{parts}
  \item Find the Fourier series expansion of the periodic function $f(x) = x^2$ in the interval $-\pi < x < \pi$. Comment on the nature of convergence of the series. Using this series, show that:
    
\begin{parts}
      \item $\sum_{n=1}^{\infty} \frac{1}{n^2} = \frac{\pi^2}{6}$
      \item $\sum_{n=1}^{\infty} \frac{(-1)^{n+1}}{n^2} = \frac{\pi^2}{12}$
    \end{parts} \pts{9}
  \item Provide one example each, with a brief explanation, of physics problems in which the following special functions and distributions occur: Legendre polynomials, Bessel functions, Hermite polynomials, Laguerre polynomials, Gamma function, and Dirac delta function. \pts{5}
\end{parts}

\vfill
\rule{\linewidth}{0.4pt}

\begin{center}\small
  \href{https://bhu-exam.vercel.app/}{Click here to visit our website}\\[4pt]
  \href{https://me-aryan.vercel.app/}{Click here to visit our contributor}\\[4pt]
  \href{https://quashan-me.vercel.app/}{Click here to visit our contributor}
\end{center}

\end{document}
